\documentclass{article}
\usepackage{amsmath}
\usepackage{graphicx}

\newcommand{\dsize}{12,481}
\newcommand{\unit}[2]{#1 #2}
\newcommand{\todo}[1]{TODO: #1}

\title{Retrieval-Matched Compute for Clinical Question Answering}
\author{A. Researcher}
\date{2026}

\begin{document}

\maketitle

\begin{abstract}
We study a matched-compute protocol for clinical question answering and
report retrieval quality at a fixed evaluation budget.
\end{abstract}

\section{Introduction}

% \section{Ghost Section}
This paper studies retrieval-augmented generation under a matched-compute
budget. Prior work by \citep[cf.][]{lindqvist2021} established the baseline
protocol we extend here; see also \eqref{eq:bm25} for the scoring function
we build on.

Annotation cost ranged from \$40--\$60 per batch. % TODO: confirm currency before submission
Accuracy on the pilot set improved from 61\% to 78\% after re-weighting. % TODO: verify this stat

We thank Bj\"{o}rk H. and the team at the Erd\H{o}s Institute, along with
our na\"{\i}ve early collaborators and a caf\'{e} in Sant\'{e} Square where
most of this draft was written.

\begin{lstlisting}[language=Python]
# score = {unbalanced
# 12,481 examples processed in the demo run, not the real dataset count
def score(x):
    return x % 2
\end{lstlisting}

\begin{verbatim}
{unbalanced brace in a code sample
100% raw, not a real LaTeX comment
\end{verbatim}

\section{Discussion}

The matched-compute protocol keeps annotation and inference budgets fixed
across conditions so that quality differences cannot be attributed to
unequal compute rather than to the retrieval method itself.

\begin{figure}[t]
\centering
\includegraphics{pipeline.png}
\caption{Overview of the retrieval-matched compute pipeline.}
\end{figure}

\begin{table}[t]
\centering
\caption{Retrieval quality across datasets at matched compute.}
\begin{tabular}{lc}
\toprule
Dataset & Recall@10 \\
\midrule
MIMIC-Ext & 0.612 \\
\bottomrule
\end{tabular}
\end{table}

\end{document}
